\documentclass[a4paper, 12pt]{report}		% general format

%%%% Charset
\usepackage{cmap}							% make PDF files searchable and copyable
\usepackage[utf8]{inputenc}					% accept different input encodings
\usepackage[T2A]{fontenc}					% russian font
\usepackage[russian]{babel}					% multilingual support (T2A)

%%%% Graphics
\usepackage[dvipsnames]{xcolor}			% driver-independent color extensions
\usepackage{graphicx}						% enhanced support for graphics
\usepackage{wrapfig}						% pro­duces fig­ures which text can flow around

%%%% Math
\usepackage{amsmath}						% Amer­i­can Math­e­mat­i­cal So­ci­ety (AMS) math fa­cil­i­ties
\usepackage{amsfonts}						% fonts from the AMS
\usepackage{amssymb}						% additional math symbols

%%%% Ty­po­grapy (don't forget about cm-super)
\usepackage{microtype}						% sublim­i­nal re­fine­ments to­wards ty­po­graph­i­cal per­fec­tion
\linespread{1.3}							% line spacing
\usepackage[left=2.5cm, right=1.5cm, top=2.5cm, bottom=2.5cm]{geometry}
\setlength{\parindent}{0pt}					% we don't want any paragraph indentation
\renewcommand{\chaptername}{}

%%%% Other
\usepackage{hyperref}							% ver­ba­tim with URL-sen­si­tive line breaks
%\DeclareUnicodeCharacter{00A0}{~}
\usepackage{float}

%------------------------------------------------------------------------------
\usepackage{listings}						% type­set source code list­ings

% Цвета для кода
\definecolor{string}{HTML}{101AF9}			% цвет строк в коде
\definecolor{comment}{HTML}{3F7F5F}		% цвет комментариев в коде
\definecolor{keyword}{HTML}{5F1441}		% цвет ключевых слов в коде
\definecolor{morecomment}{HTML}{8000FF}	% цвет include и других элементов в коде
\definecolor{captiontext}{HTML}{FFFFFF}	% цвет текста заголовка в коде
\definecolor{captionbk}{HTML}{999999}		% цвет фона заголовка в коде
\definecolor{bk}{HTML}{FFFFFF}				% цвет фона в коде
\definecolor{frame}{HTML}{999999}			% цвет рамки в коде

% Настройки отображения кода
\lstset{
	language=C++,							% Язык кода по умолчанию
	morekeywords={*,...},					% если хотите добавить ключевые слова, то добавляйте
	% Цвета
	keywordstyle=\color{keyword}\ttfamily\bfseries,
	stringstyle=\color{string}\ttfamily,
	commentstyle=\color{comment}\ttfamily\itshape,
	morecomment=[l][\color{morecomment}]{\#},
	% Настройки отображения
	breaklines=true,						% Перенос длинных строк
	basicstyle=\ttfamily\footnotesize,		% Шрифт для отображения кода
	backgroundcolor=\color{bk},				% Цвет фона кода
	%frame=lrb,xleftmargin=\fboxsep,xrightmargin=-\fboxsep, % Рамка, подогнанная к заголовку
	frame=tblr								% draw a frame at all sides of the code block
	rulecolor=\color{frame},				% Цвет рамки
	tabsize=2,								% tab space width
	showstringspaces=false,					% don't mark spaces in strings
	% Настройка отображения номеров строк. Если не нужно, то удалите весь блок
	numbers=left,							% Слева отображаются номера строк
	stepnumber=1,							% Каждую строку нумеровать
	numbersep=5pt,							% Отступ от кода
	numberstyle=\small\color{black},		% Стиль написания номеров строк
	% Для отображения русского языка
	extendedchars=true,
    literate=
        {Ö}{{\"O}}1                    {Ä}{{\"A}}1                    {Ü}{{\"U}}1
        {ß}{{\ss}}1                    {ü}{{\"u}}1                    {ä}{{\"a}}1
        {ö}{{\"o}}1                    {~}{{\textasciitilde}}1        {а}{{\selectfont\char224}}1
        {б}{{\selectfont\char225}}1    {в}{{\selectfont\char226}}1    {г}{{\selectfont\char227}}1
        {д}{{\selectfont\char228}}1    {е}{{\selectfont\char229}}1    {ё}{{\"e}}1
        {ж}{{\selectfont\char230}}1    {з}{{\selectfont\char231}}1    {и}{{\selectfont\char232}}1
        {й}{{\selectfont\char233}}1    {к}{{\selectfont\char234}}1    {л}{{\selectfont\char235}}1
        {м}{{\selectfont\char236}}1    {н}{{\selectfont\char237}}1    {о}{{\selectfont\char238}}1
        {п}{{\selectfont\char239}}1    {р}{{\selectfont\char240}}1    {с}{{\selectfont\char241}}1
        {т}{{\selectfont\char242}}1    {у}{{\selectfont\char243}}1    {ф}{{\selectfont\char244}}1
        {х}{{\selectfont\char245}}1    {ц}{{\selectfont\char246}}1    {ч}{{\selectfont\char247}}1
        {ш}{{\selectfont\char248}}1    {щ}{{\selectfont\char249}}1    {ъ}{{\selectfont\char250}}1
        {ы}{{\selectfont\char251}}1    {ь}{{\selectfont\char252}}1    {э}{{\selectfont\char253}}1
        {ю}{{\selectfont\char254}}1    {я}{{\selectfont\char255}}1    {А}{{\selectfont\char192}}1
        {Б}{{\selectfont\char193}}1    {В}{{\selectfont\char194}}1    {Г}{{\selectfont\char195}}1
        {Д}{{\selectfont\char196}}1    {Е}{{\selectfont\char197}}1    {Ё}{{\"E}}1
        {Ж}{{\selectfont\char198}}1    {З}{{\selectfont\char199}}1    {И}{{\selectfont\char200}}1
        {Й}{{\selectfont\char201}}1    {К}{{\selectfont\char202}}1    {Л}{{\selectfont\char203}}1
        {М}{{\selectfont\char204}}1    {Н}{{\selectfont\char205}}1    {О}{{\selectfont\char206}}1
        {П}{{\selectfont\char207}}1    {Р}{{\selectfont\char208}}1    {С}{{\selectfont\char209}}1
        {Т}{{\selectfont\char210}}1    {У}{{\selectfont\char211}}1    {Ф}{{\selectfont\char212}}1
        {Х}{{\selectfont\char213}}1    {Ц}{{\selectfont\char214}}1    {Ч}{{\selectfont\char215}}1
        {Ш}{{\selectfont\char216}}1    {Щ}{{\selectfont\char217}}1    {Ъ}{{\selectfont\char218}}1
        {Ы}{{\selectfont\char219}}1    {Ь}{{\selectfont\char220}}1    {Э}{{\selectfont\char221}}1
        {Ю}{{\selectfont\char222}}1    {Я}{{\selectfont\char223}}1    {і}{{\selectfont\char105}}1
        {ї}{{\selectfont\char168}}1    {є}{{\selectfont\char185}}1    {ґ}{{\selectfont\char160}}1
        {І}{{\selectfont\char73}}1     {Ї}{{\selectfont\char136}}1    {Є}{{\selectfont\char153}}1
        {Ґ}{{\selectfont\char128}}1
}

% Для настройки заголовка кода
\usepackage{caption}
\DeclareCaptionFont{white}{\color{сaptiontext}}
\DeclareCaptionFormat{listing}{\parbox{\linewidth}{\colorbox{сaptionbk}{\parbox{\linewidth}{#1#2#3}}\vskip-4pt}}
%\captionsetup[lstlisting]{format=listing,labelfont=white,textfont=white}
\renewcommand{\lstlistingname}{Листинг} % Переименование Listings в нужное именование структуры

%------------------------------------------------------------------------------
\begin{document}

\begin{titlepage}

%----------------------------------------------------------------------------------------
%	HEADING SECTIONS
%----------------------------------------------------------------------------------------
\begin{center} % Center everything
Федеральное государственное автономное образовательное \\
учреждение высшего образования \\[0.4cm]
% https://www.spbstu.ru/university/organizational-documents/corporate-identity/
\includegraphics[scale=0.8]{res/SPbPU-logo} \\[0.4cm]
Институт компьютерных наук и технологий \\*
Высшая школа интеллектуальных систем и суперкомпьютерных технологий
\end{center}

\vspace{3cm}

%----------------------------------------------------------------------------------------
%	TITLE SECTION
%----------------------------------------------------------------------------------------
\begin{center} % Center everything
\textbf{Курсовая работа}\\
по дисциплине "Сетевая безопасность"\\
на тему\\
\textbf{Исследование коммуникационного протокола WebSocket}
\end{center}

\vspace{3.5cm}
 
%----------------------------------------------------------------------------------------
%	AUTHOR SECTION
%----------------------------------------------------------------------------------------
\begin{flushleft}
Выполнил студент гр. 3540901/21501 \hspace{3cm} $\underset{\text{(подпись)}}{\underline{\hspace{3cm}}}$ С.А.Мартынов\\[0.5cm]
Преподаватель \hspace{7.25cm} $\underset{\text{(подпись)}}{\underline{\hspace{3cm}}}$ В.Э. Шмаков\\[0.5cm]
\hspace{10.2cm} «\underline{\hspace{1cm}}» \underline{\hspace{3cm}} 2023 г.
\end{flushleft}

\vfill % Fill the rest of the page with whitespace

%----------------------------------------------------------------------------------------
%	DATE SECTION
%----------------------------------------------------------------------------------------
\begin{center}
Санкт-Петербург\\
2023
\end{center}

\end{titlepage}

\setcounter{page}{2}                            % inclide the title page
\tableofcontents
%\input{text}                                 % inclide the main text
%% \newpage
\section*{}
\addcontentsline{toc}{section}{List of sources}

\begin{thebibliography}{00}

% Use: \cite{label1}
\bibitem{label1} JavaScript.ru: Справочник по работе с WebSocket. URL: https://learn.javascript.ru/websockets (дата обращения: 16.04.2023).

\bibitem{label2} Internet Engineering Task Force (IETF) / Request for Comments: 6455: The WebSocket Protocol (December 2011). URL: https://datatracker.ietf.org/doc/html/rfc6455 (дата обращения: 20.04.2023).

\bibitem{label3} MDN Web Docs / The WebSocket API (WebSockets). URL: https://developer.mozilla.org/en-US/docs/Web/API/WebSockets\_API (дата обращения: 20.04.2023).

\bibitem{label4} Nginx blog / NGINX as a WebSocket Proxy (Rick Nelson, 2014). URL: https://www.nginx.com/blog/websocket-nginx/ (дата обращения: 23.04.2023).

\bibitem{label5} How to scale WebSocket – horizontal scaling with WebSocket tutorial (2020). URL: https://tsh.io/blog/how-to-scale-websocket/ (дата обращения: 26.04.2023).

\end{thebibliography}                              % inclide the list of sources


\end{document}

%------------------------------------------------------------------------------
% Examples:
%
%
%
% \begin{figure}[H]
% \centering
% \includegraphics[scale=0.8]{res/pic01}
% \caption{Picture description}
% \end{figure}
%
%
%
% \begin{table}[htb]
%     \begin{tabularx}{\textwidth}{|X|c|c|c|c|c|}
%     \hline
%     \multirow{2}{*}{tb1} & \multirow{2}{*}{tbl2} & \multicolumn{4}{c|}{tbl3} \\
%     \cline{3-6}
%     {} & {} & A & B & C & D\\
%     \hline
%     Text & {} & Text & {} & {} & {} \\
%     \hline
%     \end{tabularx}
% \caption{Table description}
% \end{table}